\input{../header}

\rhead{Your name: \rule{8cm}{0.15mm}}

\begin{document}

\section{Learning target I1, version 1}
Compute approximations of definite integrals using numerical methods. Decide when each method gives overestimates vs. underestimates.

\hrulefill

Consider the function \(f(x)=\dfrac{\sin(x)}{x}\), which does not have an elementary antiderivative. Define the integral function \(F(t) = \displaystyle\int_1^t f(x)\ dx\).
\begin{enumerate}[(a)]
  \item What does the Second FTC say about the relationship between $F$ and $f$?
  \vfill

  \item Use appropriate technology, such as Geogebra or a Riemann sum applet, to approximate $F(2)$ using left, right, midpoint, and trapezoid sums with 10 subintervals:
  \begin{itemize}
    \item $L_{10}=$
    \item $R_{10}=$
    \item $M_{10}=$
    \item $T_{10}=$
  \end{itemize}
  Report your answers with at least 6 decimal places of accuracy.
  
  \item Which of these approximations are overestimates and which are underestimates? \\
  How can you tell?
  \vfill

  \item Use your calculations from part (b) to approximate $F(2)$ using Simpson's rule.
\end{enumerate}
\end{document}